\section{HMM 把隐变量排列成时间链}
\label{sec:13-Machine-Learning/06-Probabilistic-Models/05}

你有个朋友每天发一条消息，只写两个字：开心，或者不开心。你翻了三十天的记录，发现开心和不开心有某种黏性——连续几天开心之后很少猛地蹦出个不开心，反过来也一样。你猜背后有个你看不见的东西在操纵这一切：也许是天气。晴天心情好的概率大，雨天心情差的概率大。但你看不到天气，只能看到每天那两个字的记录。

你想从三十天的记录出发做几件事。第一，估算每一天背后是晴还是雨的概率。第二，算出整段三十天最可能的那一条天气序列。第三，从数据里把转移概率和发射概率学出来。

GMM 假设每个样本的隐类别彼此独立——今天的状态和昨天没有任何关系。语音音素、设备工况、说话人的情绪阶段却天然有持续性。Hidden Markov model 做的事就是把隐变量排成一条时间链：

\[
p(z_{1:T},x_{1:T})
=p(z_1)
\prod_{t=2}^Tp(z_t\mid z_{t-1})
\prod_{t=1}^Tp(x_t\mid z_t).
\]

两颗钉子固定了整个模型。第一颗：隐状态是一阶 Markov——\(z_t\) 只依赖 \(z_{t-1}\)，再往前的历史通过 \(z_{t-1}\) 已经汇总完毕。第二颗：给定当前隐状态后，当前观测和过去未来所有其他变量条件独立——\(x_t\) 只看 \(z_t\) 的脸色，不直接看 \(z_{t-1}\) 或 \(x_{t+1}\)。

\subsection{尝试枚举然后撞上指数墙}

最朴素的办法：有 \(K\) 种天气，\(T\) 天，枚举所有 \(K^T\) 条天气路径，每条算一个联合概率。\(K=2,T=30\) 时是 \(2^{30}\approx10^9\) 条路径——对三十个数据点要跑十亿次计算。稍微加长到一百天，\(2^{100}\) 已经超过了宇宙中可观测原子的数量级。

你跑不动不是因为你写的代码不够优化。是问题规模随序列长度指数增长，没有算法能靠枚举所有路径来解决问题。只能找办法在每一步把前缀信息压缩，不再回头展开全部历史。

\subsection{Forward：把每条路径的负担变成每个状态的汇总}

令初始分布为 \(\pi_j\)——第一天是晴是雨的先验概率。从状态 \(i\) 转移到状态 \(j\) 的概率为 \(A_{ij}\)，在状态 \(j\) 下发出观测 \(x\) 的概率密度为 \(b_j(x)\)。

定义 forward message：

\[
\alpha_t(j)=p(x_{1:t},z_t=j).
\]

读作：到时刻 \(t\) 为止已看到 \(x_1\) 到 \(x_t\)，且第 \(t\) 时刻状态恰为 \(j\) 的联合概率。它是前缀信息的一种压缩形式。

递推只有一行：

\[
\alpha_1(j)=\pi_jb_j(x_1),
\]

\[
\alpha_t(j)
=b_j(x_t)\sum_i\alpha_{t-1}(i)A_{ij}.
\]

理解这行的方式不是死记符号。时刻 \(t\) 到达状态 \(j\) 的所有路径，可以按其前一时刻的状态 \(i\) 分类。到达 \(i\) 的所有前缀概率已经被 \(\alpha_{t-1}(i)\) 汇总了；从 \(i\) 跨一步到 \(j\) 需要乘上转移概率 \(A_{ij}\)；然后在状态 \(j\) 下产生当前观测 \(x_t\) 的代价是 \(b_j(x_t)\)。对 \(i\) 求和，就是把以不同状态为结尾的所有前缀路径合并到同一个 \(j\) 上。

\[
p(x_{1:T})=\sum_j\alpha_T(j).
\]

总观测概率就是最后时刻各个状态下 forward message 的简单求和。

暴力枚举有 \(K^T\) 条路径。forward 在每个时刻只维护 \(K\) 个数字——每个可能状态一个——每个时刻的计算量是 \(O(K^2)\)。总复杂度 \(O(TK^2)\)，和序列长度线性。这就是动态规划：当前只需要知道每个状态的汇总概率，不需要知道这个概率是由哪些具体历史路径拼出来的。

\subsection{Backward 让未来信息回流修正过去}

知道整段序列后再回头看中间某个时刻的状态，和你一路往前走只用到当前为止信息时的判断，可以是完全不同的。

假设你翻完一个朋友三十天的心情记录，发现最后三天全是开心。往回推，倒数第四天大概率也是开心，因为从开心跳到不开心再从那里再跳回开心，在转移矩阵不太极端时是一条低概率的曲折路径。但如果你只看到倒数第四天为止的信息，你还不知道后面三天的走势，你当时的即时判断可能偏向中性。看到末尾再回头看，中间的模糊段被后续信息修正了。

定义为 backward message：

\[
\beta_t(i)=p(x_{t+1:T}\mid z_t=i).
\]

读作：给定当前状态为 \(i\)，未来观测从 \(t+1\) 到 \(T\) 的条件概率。从末端 \(\beta_T(i)=1\) 出发反向递推：

\[
\beta_t(i)
=\sum_j A_{ij}b_j(x_{t+1})\beta_{t+1}(j).
\]

到达 \(z_t=i\) 后，下一步跳到哪个 \(j\)（由 \(A_{ij}\) 决定），在 \(j\) 上产生观测 \(x_{t+1}\)（由 \(b_j(x_{t+1})\) 决定），然后剩下的从 \(t+1\) 往后的未来概率已经由 \(\beta_{t+1}(j)\) 汇总。

关键一步：在给定 \(z_t\) 的条件下，历史 \(x_{1:t}\) 和未来 \(x_{t+1:T}\) 条件独立。因此

\[
\alpha_t(j)\beta_t(j)=p(z_t=j,x_{1:T}),
\]

归一化即得平滑后验：

\[
p(z_t=j\mid x_{1:T})
=\frac{\alpha_t(j)\beta_t(j)}
{\sum_{j'}\alpha_t(j')\beta_t(j')}.
\]

filtering 只用 \(x_{1:t}\) 推断当前状态——适合在线的实时追踪。smoothing 用完整序列 \(x_{1:T}\) 回溯修正每个时刻——适合离线分析。两者的信息集不同，不能随口混成一件事。

\subsection{Viterbi 回答的问题和逐点最优不是同一个}

要知道整条最可能的天气序列——联合概率最大的那条——不是按时刻各自挑最可能的状态就能拼出来的。

逐时刻选 marginal posterior 最大的状态：\(t=3\) 是晴天因为概率最高，\(t=4\) 是雨天因为概率最高。但这样拼出来的序列里，一步晴天一步雨天，转移概率的乘积可能极小，甚至对应一条转移概率为零的不允许路径。每个时刻拿的是边缘最优，没有人检查相邻两帧之间是否兼容。

Viterbi 把 forward 递推里的求和换为最大化：

\[
\max_{z_{1:T}}p(z_{1:T}\mid x_{1:T}),
\]

并保存每步是从哪个前驱状态取到的最大值（回溯指针）。递推时在每个状态记录``到达我的最佳路径来自哪个前驱''，最后在末端选最大者，顺着指针倒回去就得到了完整的最优序列。

\subsection{Baum--Welch：状态不可见时两边轮流猜}

如果每条记录都标注了当天的天气，事情就简单了：转移矩阵靠数出每种转移的发生次数来估计，发射概率靠按状态分组求观测的统计量。问题是你看不到天气。

Baum--Welch 做的事情就是两边轮流。E 步：用当前的参数估计，跑 forward--backward 算出每个时刻的状态后验概率 \(p(z_t=j\mid x_{1:T})\)，以及相邻时刻的联合后验 \(p(z_t=i,z_{t+1}=j\mid x_{1:T})\)。M 步：用这些后验概率作为``软计数''——不是某条路径确切经过状态 \(j\) 的0/1指示，而是以概率加权——重新估计初始分布、转移矩阵和发射参数。然后回到 E 步，用新参数再跑一遍。

这就是 EM 算法在序列模型上的形态。它继承了 EM 的全部脾气：likelihood 单调不降，但不保证到全局最优；不同初值可能收敛到不同局部模式；状态标签可以任意置换——把状态1全改成状态2同时换掉对应的发射参数，likelihood 完全相同。算法告诉你数据里似乎有两类状态，但不会告诉你哪类对应晴天、哪类对应雨天。学出来的状态标签，需要用发射参数的物理含义、持续时间分布、外部标注和领域知识一起解读。

\subsection{数值上的暗坑：概率连乘变零}

长序列里许多接近0到1之间的小数连乘，标准浮点数会迅速下溢为零。\(\alpha_t\) 和 \(\beta_t\) 都会遇到。可以在每个时刻对 \(\alpha_t\) 做缩放——除以当前所有状态的概率和——让数值保持在合理范围，同时把缩放常数单独累计，最后用来恢复 log-likelihood。也可以在 log 域用 log-sum-exp 把所有计算留在对数空间。

一句实操建议：先在很短序列（三五个时刻）上写出所有状态路径的枚举结果，把 forward likelihood、smoothing posterior、Viterbi 路径全部对一遍，确认数值完全吻合无误，再扩展到实用长度。

\subsection{标准 HMM 认为你在每个状态的停留时间是无记忆的}

Markov 转移意味着每一步都以固定概率离开当前状态，和已经呆了多久无关。状态持续时间隐含着几何分布——短停留最常见，长停留的概率指数衰减。如果你的雨天真实有某种最短持续时间（比如一场雨很少少于半天），标准 HMM 无法自然地表达这一约束。

用更多状态去补：雨天被拆成``雨刚开始''``雨中段''``雨快停了''，用若干个状态的链来拼出一个非几何的停留时长分布。但它只是用多个状态强行拼接，并不代表真的发现了多个不同的物理机制。Hidden semi-Markov model 等扩展直接在状态上绑定停留时间分布，避免了这种状态的虚假膨胀。更一般地，如果转移概率自身随时间或外部协变量变化，需要非齐次模型。

\subsubsection{一个两状态三时刻的手算例子}

取一个两状态、三时刻的 HMM，枚举全部八条隐状态路径，filtering、smoothing 与 Viterbi 路径都可以手算到底。三者的对象不同：filtering 只用到当前时刻为止的观测，smoothing 用整段序列修正每个时刻，Viterbi 返回联合概率最大的一条完整路径。

把转移矩阵调得很粘——留在原状态的概率接近一——可以找到这样的参数设置：某时刻 Viterbi 路径经过的状态并不是该时刻 marginal posterior 最大的状态。这一点值得亲手算一遍，算过之后就不会再把逐点最优和全局最优混为一谈。

把序列复制到很长，普通概率域的连乘会下溢为零；缩放版与 log 域版则给出一致且有限的 log-likelihood。在手算例子中把缩放系数的变化也走一遍，以后在代码里看到概率突然变成零就不会慌张。

随机过程层面的基础可回看\hyperref[sec:06-Random-Processes/04-Applications/04]{``Hidden Markov Model：从观测反推状态''}。
